\subsection{The case of solvable groups}

Let \(\mathbf{E}\) be a real vector space, and \(\mathbf{K}\) a convex subset of \(\mathbf{E}\) . A mapping \(u: \mathbf{K} \to \mathbf{K}\) such that for \(x, y\) in \(\mathbf{K}\) and for every real number \(t\) in \([0, 1]\) , we have

\[
u (t x + (1 - t) y) = t u (x) + (1 - t) u (y)\tag{1}
\]

is said to be an affine transformation. From relation (1) we deduce that

\[
u \left(\sum_ {i \in \mathrm{I}} t _ {i} x _ {i}\right) = \sum_ {i \in \mathrm{I}} t _ {i} u \left(x _ {i}\right)\tag{2}
\]

for every finite set I, points \(x_{i}\) in K and positive real numbers \(t_i\) such that \(\sum_{i\in I}t_i = 1\) .

Let u and v be two affine transformations on K, then the mapping \(u \circ v\) is an affine transformation on K. If \(v: E \to E\) is a linear mapping such that \(v(\mathbf{K}) \subset \mathbf{K}\) , the mapping \(u: K \to K\) which coincides with v on K is an affine transformation.

THEOREM 1 (Markoff-Kakutani). --- Let E be a Hausdorff locally convex vector space over the field R, and K a non-empty compact convex subset of E. Let \(\Gamma\) be a set of continuous affine transformations on K, pairwise permutable. Then there exists a point a in K such that \(u(a) = a\) for all \(u \in \Gamma\) .

For every \(u \in \Gamma\) , let \(\mathbf{K}_u\) be the set of all \(x \in \mathbf{K}\) such that \(u(x) = x\) . We shall show that \(\mathbf{K}_u\) is non-empty. Let \(x\) be a point of \(\mathbf{K}\) ; for every integer \(n \geqslant 1\) , let \(x_n\) be the element \(\frac{1}{n} \sum_{i=0}^{n-1} u^i(x)\) of \(\mathbf{E}\) . Since \(\mathbf{K}\) is convex and stable under \(u\) , the points \(x_n\) belong to \(\mathbf{K}\) and since \(\mathbf{K}\) is compact, there exists a limit point \(a\) of the sequence \((x_n)_{n \geqslant 1}\) . The mapping \(y \mapsto u(y) - y\) from \(\mathbf{K}\) into \(\mathbf{E}\) is continuous, hence \(u(a) - a\) is a limit point of the sequence \((u(x_n) - x_n)_{n \geqslant 1}\) . But we have \(u(x_n) - x_n = \frac{1}{n} (u^n(x) - x)\) .

Since K is compact, hence also bounded (III, p.~3, prop. 2), the sequence \((u^n(x) - x)_{n\geqslant 1}\) is bounded; consequently, the sequence \(\left(\frac{1}{n} (u^n (x) - x)\right)_{n\geqslant 1}\) tends to 0 (III, p.~4, prop. 3), and since E is Hausdorff, we have \(u(a) - a = 0\) . Therefore \(a\in \mathbf{K}_u\) .

Each of the sets \(\mathbf{K}_u\) is a closed and convex subset of the compact space \(\mathbf{K}\) , and we shall prove that the intersection \(\bigcap_{u\in \Gamma}\mathbf{K}_u\) is non-empty. For this it is enough to

prove that, for \(n \geqslant 1\) , and \(u_{1}, \ldots, u_{n}\) in \(\Gamma\) , the set \(K_{u_{1}} \cap \ldots \cap K_{u_{n}}\) is non-empty. The case n = 1 having been considered, we argue by induction on n.~Suppose \(n \geqslant 2\) and put \(L = K_{u_{1}} \cap \ldots \cap K_{u_{n-1}}\) . By the hypothesis of induction, L is a non-empty compact convex subset of E. Since \(u_{n}\) commutes with \(u_{1}, \ldots, u_{n-1}\) , we have \(u_{n}(L) \subset L\) . Applying the first part of the proof to the affine transformation induced by \(u_{n}\) on L, we conclude that there exists a point a in L such that \(u_{n}(a) = a\) ; then a belongs to \(K_{u_{1}} \cap \ldots \cap K_{u_{n}}\) , which is then non-empty.

COROLLARY. --- Let G be a solvable group of continuous affine transformations on K. Then there exists a point in K which is invariant under G.

By the definition of a solvable group (A, I, § 6, No.~4) there exists a finite decreasing sequence \((\mathbf{G}_i)_{0\leqslant i\leqslant n}\) of distinct subgroups of \(\mathbf{G}\) , such that \(\mathbf{G}_0 = \mathbf{G}\) , \(\mathbf{G}_n = \{e\}\) and such that the group \(\mathbf{G}_{i - 1} / \mathbf{G}_i\) is commutative for \(1\leqslant i\leqslant n\) . Let \(\mathbf{K}_i\) denote the set of fixed points of \(\mathbf{G}_i\) in \(\mathbf{K}\) . Then \(\mathbf{K}_n = \mathbf{K}\) . Moreover, for \(1\leqslant i\leqslant n\) , every element of \(\mathbf{G}_i\) induces the identity transformation on \(\mathbf{K}_i\) ; we thus deduce an action of the abelian group \(\mathbf{G}_{i - 1} / \mathbf{G}_i\) on \(\mathbf{K}\) if \(\mathbf{K}_i\) is non-empty; it follows from th. 1 that the set \(\mathbf{K}_{i - 1}\) of fixed points of \(\mathbf{G}_{i - 1} / \mathbf{G}_i\) in \(\mathbf{K}_i\) is non-empty. By descending induction on \(i\) , we conclude that \(\mathbf{K}_0\) is non-empty, hence the corollary.

